% This is main_kpk_v03.tex - Master Academic Manuscript for Springer LNCS / ICCSCI (Procedia Computer Science)
% Full Comprehensive & Non-Summarized Version (Zero Data Loss)
\documentclass[runningheads]{llncs}

% Core Packages
\usepackage{graphicx}
\usepackage{amsmath,amssymb,amsfonts}
\usepackage{booktabs}
\usepackage{array}
\usepackage{adjustbox}
\usepackage{algorithm}
\usepackage{algorithmic}
\usepackage{url}
\usepackage{hyperref}
\usepackage{xcolor}

% Configure graphics search path for local execution and Overleaf environments
\graphicspath{{charts/v03/}{./charts/v03/}{../../charts/v03/}{src/v3-run/}}

\begin{document}

\title{Unsupervised Dual-Path Consensus Machine Learning and Operational Policy Framework for Village Fund Expenditure Anomaly Detection}

\titlerunning{Dual-Path Consensus Anomaly Detection for Village Funds}

\author{Research Team\inst{1}\orcidID{0000-0000-0000-0000}}
\authorrunning{Research Team et al.}

\institute{Faculty of Computer Science, University Locus, Indonesia\\
\email{author@domain.ac.id}}

\maketitle

\begin{abstract}
Indonesia's \textit{Dana Desa} (Village Fund) channels Rp 71 trillion annually to 75,259 villages, yet legal prosecutions reveal systemic financial diversion lagging two to five years behind disbursement windows. Existing monitoring relies on retrospective manual audits or supervised machine learning requiring non-existent ground-truth fraud labels in real-time administrative \textit{Siskeudes} records. Grounded in Design Science Research (DSR), Agency Theory, and Fraud Diamond, this study constructs an unsupervised Dual-Path Consensus framework (\textbf{Protocol 2}) combining Isolation Forest (IF), Local Outlier Factor (LOF), and an 8-layer Reconstruction Dense Autoencoder (RDA) with an Operational Policy Mapping Layer on 96,778 activity records across 1,363 villages (2023--2025) in Jambi Province. Compared to baseline majority voting (\textbf{Protocol 1}), \textbf{Protocol 2} decouples local density isolates (LOF) from multi-model global convergence (IF $\cap$ RDA), expanding anomaly recall from 3,107 (3.12\%) to \textbf{7,153 consensus anomalous activities (7.39\%)} representing \textbf{Rp 642.85 billion in financial realization exposure (14.89\%)} and \textbf{Rp 6.08 trillion in budget ceiling exposure (7.46\%)}. Operational policy mapping reveals a major typology shift toward \textbf{T2: Ghost Activities (58.1\% / 4,155 records)} and \textbf{T5: Procurement Irregularities (32.8\% / 2,343 records)}. Longitudinal persistence identifies \textbf{702 villages (51.5\%)} flagged across all three fiscal years. On an ex-ante synthetic benchmark, \textbf{Protocol 2} achieves superior recovery (\textbf{Precision = 0.846, Recall = 0.846, F1 = 0.846, AUC-ROC = 0.912} vs. \textbf{Protocol 1} F1 = 0.612). Operationalising the DeLone \& McLean IS Success Model, \textbf{Protocol 2} reduces inspectorate search space by 92.6\% and translates neural loss into Explainable AI (XAI) audit checklists for APIP field inspections.

\keywords{Village Fund Governance \and Unsupervised Anomaly Detection \and Design Science Research \and Dual-Path Consensus \and Explainable AI \and Corruption Typology.}
\end{abstract}

% -------------------------------------------------------------------------
% SECTION 1: INTRODUCTION
% -------------------------------------------------------------------------
\section{Introduction}

Indonesia's \textit{Dana Desa} (Village Fund) programme represents one of the largest intergovernmental fiscal transfers in the developing world, channelling approximately Rp 71 trillion annually to 75,259 villages across the archipelago since the enactment of Law No. 6 of 2014 on Villages~\cite{ref1}. The programme's fiscal scale, however, has severely outpaced the supervisory infrastructure needed to ensure its financial integrity. By 2024, Indonesian Corruption Watch (ICW) had catalogued 591 court verdicts involving village fund misappropriation, implicating 640 defendants and documenting Rp 598.13 billion in verified state financial losses~\cite{ref2}. Concurrently, the Komisi Pemberantasan Korupsi (KPK) identified 851 village fund corruption cases from 2015 onward, in which village heads (\textit{Kepala Desa}) accounted for over 60\% of perpetrators~\cite{ref3}. These empirical figures converge on a single diagnostic conclusion: the dominant institutional response to village fund corruption remains reactive---prosecution occurring two to five years after fund diversion---rather than proactive, data-driven detection operating within active disbursement cycles.

The technical infrastructure for proactive detection already exists in latent administrative databases. The \textit{Sistem Keuangan Desa} (Siskeudes), maintained by BPKP and the Ministry of Home Affairs (\textit{Kemendagri}), records village fund expenditure at the granular activity level across each disbursement stage: what type of output activity was planned, how much budget was allocated (\textit{Pagu}), how much was realised per tranche (T1, T2, T3), and through what procurement mechanism (\textit{Cara Pengadaan}). This granular absorption data constitutes an administrative signal archive---one that encodes, in latent statistical form, the expenditure patterns that distinguish legitimate governance from fraudulent diversion. Yet no systematic, automated screening mechanism currently operates at the district or provincial level to interrogate this archive before inspection cycles conclude and disbursement windows close.

Existing research on \textit{Dana Desa} fraud concentrates overwhelmingly on retrospective forensic analysis~\cite{ref4}, qualitative governance accountability frameworks~\cite{ref5}, and traditional fraud triangle diagnostics~\cite{ref6}. Where machine learning enters the literature, it applies predominantly supervised classification---a paradigm that requires pre-existing fraud labels unavailable in real-time administrative monitoring contexts~\cite{ref13,ref14}. The present study addresses this gap directly through a \textbf{Design Science Research (DSR)} approach~\cite{ref5,ref10}: it designs, constructs, and evaluates a multi-paradigm unsupervised learning artifact (\textbf{Protocol 2}) that detects expenditure anomalies across a longitudinal panel of 96,778 activity-level expenditure records (1,363 villages, 3 fiscal years 2023--2025) in Jambi Province, without presupposing any prior fraud labels.

\textbf{Jambi Province provides a structurally significant research site.} The \texttt{jaga.id} civil-society transparency platform, operated by KPK, recorded only 11 complaint submissions attributed to Jambi Province---1.4\% of 761 national reports---while provinces of comparable administrative scale registered substantially higher reporting volumes~\cite{ref26}. Low voluntary public reporting does not signal low corruption exposure; rather, it signals a low-monitoring environment. Alfada's panel GMM analysis across 19 Indonesian provinces confirms that regions combining high intergovernmental transfer dependence with weak accountability structures exhibit systematically elevated corruption incidence---the structural profile directly applicable to Jambi~\cite{ref27}. Furthermore, district-level \textit{Aparat Pengawasan Intern Pemerintah} (APIP) inspectorates face severe operational capacity constraints: Srirejeki and Faturokhman~\cite{ref12} show that kabupaten APIP offices lack staffing (typically only 5 to 15 auditors per kabupaten) to manually screen thousands of individual village activity records per disbursement cycle, producing a measurable shortfall between the statutory monitoring mandate under Government Regulation No. 60/2008 and actual physical inspection coverage.

Contemporaneous court verdicts in Jambi Province confirm that detection lags range from two to five years across jurisdictions:
\begin{enumerate}
    \item \textit{Kasus Muara Hemat (Kerinci)}: Verified state loss of Rp 644 million, with an enforcement lag of 5 years~\cite{ref28,ref30}.
    \item \textit{Kasus Jambi Tulo (Muaro Jambi)}: Village fund freezing due to fictitious project indications exceeding Rp 300 million~\cite{ref29}.
    \item \textit{Kasus Mantan Kades Kerinci}: Misappropriation of physical construction funds amounting to Rp 644 million~\cite{ref30}.
    \item \textit{Kasus Pangkal Duri (Tanjung Jabung Timur)}: State financial loss reaching Rp 415 million~\cite{ref31}.
\end{enumerate}
All four judicial cases involved financial irregularities that generated quantifiable signals in Siskeudes absorption records: fictitious construction outputs, zero-realisation procurement activities, and uncompetitive Swakelola budget execution---precisely the patterns targeted by the feature engineering and Dual-Path consensus architecture in this study.

This study makes three main contributions:
\begin{enumerate}
    \item \textbf{Methodological Contribution (DSR Artifact)}: It constructs an unsupervised Dual-Path Consensus Machine Learning framework combining local density estimation (LOF) and global anomaly convergence (IF $\cap$ RDA) with an Operational Policy Mapping Layer, resolving the mutual cancellation limitation of traditional majority voting.
    \item \textbf{Theoretical Contribution}: It operationalises Agency Theory, the Fraud Diamond, and Principal-Agent information asymmetry within administrative Siskeudes records, bridging Information Systems (IS) governance literature and computational fraud detection.
    \item \textbf{Practical Contribution}: It provides an ex-ante evaluation framework that reduces APIP inspection search space by 92.6\% at the activity level and prioritises 702 persistent high-risk villages (Tier 1), generating instance-level Explainable AI (XAI) audit checklists for field inspectors.
\end{enumerate}

This study addresses three core Research Questions (RQs):
\begin{itemize}
    \item \textbf{RQ1}: Which feature constructs derived from Siskeudes expenditure absorption data serve as the most discriminating signals of expenditure anomaly?
    \item \textbf{RQ2}: Which algorithmic paradigm---Isolation Forest, Local Outlier Factor, or Reconstruction Dense Autoencoder---and ensemble gating mechanism demonstrates superior anomaly discrimination on longitudinal village fund data?
    \item \textbf{RQ3}: How do algorithmically identified anomalies map to established corruption typologies, and how can XAI loss decomposition translate mathematical outliers into actionable APIP inspection checklists?
\end{itemize}

% -------------------------------------------------------------------------
% SECTION 2: RELATED WORK
% -------------------------------------------------------------------------
\section{Related Work}

\subsection{Corruption Typology Frameworks}
Understanding the target of unsupervised detection requires a precise, domain-grounded taxonomy of corruption forms. Bussell~\cite{ref4} proposes a pragmatic two-dimensional typology distinguishing corruption by access type (monetary vs. preferential) and governance level (bureaucratic vs. political). Graycar's TASP framework---Types, Activities, Sectors, Places---extends this by asserting that anti-corruption mechanisms must specify not merely what form of corruption occurs, but within which specific public programme activity and geographic locus it manifests~\cite{ref7}.

For Indonesia's \textit{Dana Desa} context, Siregar and Aminudin~\cite{ref13} provide the primary empirical classification: a multi-case study identifying five recurring modus operandi---unit price mark-up of goods and services, fictitious budget items (\textit{proyek fiktif}), double budgeting, procurement manipulation, and misuse of personnel honoraria. Kartadinata et al.~\cite{ref14} extend this taxonomy through systematic analysis of over 200 KPK prosecution cases, confirming mark-up (\textit{penggelembungan}) and fictitious projects as the dominant physical corruption moduses. Recent empirical work by Medan et al.~\cite{ref15} documents that these exact modus operandi patterns persist into contemporary fiscal periods, validating typology-grounded detection rules.

\subsection{Fraud Diamond and Principal-Agent Foundations}
The Fraud Triangle (Cressey, 1953)~\cite{ref17}, operationalised in rural governance by Hidajat~\cite{ref6}, models fraud via three conditions: Pressure (rigid statutory disbursement timetables), Opportunity (structural absence of competitive bidding), and Rationalisation. Wolfe and Hermanson's Fraud Diamond extends this model by incorporating \textbf{Capability}---the individual authority, digital access, and technical competence required to execute and conceal financial manipulation~\cite{ref17b}. In Siskeudes execution, Capability resides in the Village Head (\textit{Kepala Desa}) and Financial Officer (\textit{Kaur Keuangan}), who maintain exclusive digital authorization credentials and bank withdrawal specimen signatures.

Opportunity is structurally driven by procurement method selection. Dataset audit reveals that \textbf{98.8\% of all village fund activities} in Jambi Province are executed through self-managed procurement (\textit{Swakelola}), completely bypassing open competitive bidding. This empirical dominance confirms S{\o}reide's~\cite{ref9} procurement corruption theory: the structural absence of market competition removes natural price-discovery mechanisms, enabling local officials to select favored vendors or inflate invoices without external challenge.

Principal-Agent Theory~\cite{ref25} frames the village head as an Agent possessing severe information asymmetry relative to the Principal (district APIP inspectorate, BPKP, KPK):
\begin{equation}
\text{InfoAsymmetry} = \mathcal{I}_{\text{Agent}}(\text{PhysicalRealization}, \text{TrueCost}) - \mathcal{I}_{\text{Principal}}(\text{SiskeudesReports})
\end{equation}
Unexplained statistical deviations in unit costs, tranche disbursement progress ratios, and procurement categories constitute proximate empirical traces of moral hazard and information asymmetry exploitation~\cite{ref9,ref25}.

\subsection{Anomaly Detection Paradigms in Public Financial Data}
The anomaly detection literature applied to public financial management has converged on three methodologically distinct paradigms:
\begin{enumerate}
    \item \textbf{Global Sparsity Partitioning}: Isolation Forest (Liu et al.~\cite{ref18,ref19}) measures random partitioning path length in tree ensembles, performing optimally against globally extreme multi-feature outliers.
    \item \textbf{Local Density Ratio Estimation}: Local Outlier Factor (LOF, Breunig et al.~\cite{ref24}) evaluates local reachability density ratios relative to $k$-nearest neighbors, isolating records that are locally sparse relative to their specific activity peer group (e.g., within a specific output code \texttt{Kode\_Output}).
    \item \textbf{Neural Reconstruction Error}: Deep Autoencoders (Zhou \& Paffenroth~\cite{ref32}) learn compressed latent representations of normal financial behaviour. Records exhibiting anomalously high Mean Squared Error (MSE) at reconstruction signal non-linear feature interaction distortions.
\end{enumerate}

\subsection{Information Systems Success Grounding}
The DeLone and McLean IS Success Model~\cite{ref10}---in its 6-box flow incorporating System Quality $\to$ Information Quality $\to$ Use $\to$ Individual Impact $\to$ Organizational Impact---provides the theoretical justification for deploying an automated anomaly detection artifact. In this framework, processed outputs (dual-path consensus flags, typology mappings, village priority tiers) represent \textbf{Information Quality}. High Information Quality directly drives \textbf{Individual Impact} (reducing auditor search space by 92.6\% and providing XAI checklists) and \textbf{Organizational Impact} (enabling targeted APIP inspection and state financial loss recovery). Mutungi et al.~\cite{ref5} demonstrate that digital anti-corruption tools fail when algorithmic design does not map directly to specific administrative interaction points.

% -------------------------------------------------------------------------
% SECTION 3: METHODOLOGY
% -------------------------------------------------------------------------
\section{Methodology}

\subsection{Design Science Research (DSR) Framework}
This study adheres strictly to the Design Science Research (DSR) methodology for Information Systems~\cite{ref5,ref10}, structuring artifact iteration across three primary cycles:
\begin{enumerate}
    \item \textbf{Relevance Cycle}: Addresses the institutional monitoring problem across 96,778 village fund activity entries under severe APIP district auditor capacity constraints.
    \item \textbf{Rigor Cycle}: Grounds feature construction in Agency Theory and Fraud Diamond dynamics, drawing algorithms from multi-paradigm anomaly detection literature~\cite{ref18,ref24,ref32}.
    \item \textbf{Design Cycle}: Iteratively constructs the IT artifact comprising robust feature engineering, a Dual-Path Consensus Machine Learning engine, and an Operational Policy Mapping Layer.
\end{enumerate}

\subsection{Dataset Cleaning \& Record Filtering Audit}
The dataset combines two administrative data sources collected via the KPK \texttt{jaga.id} portal~\cite{ref26}: (1) \textbf{Penyerapan} (expenditure absorption) records, and (2) \textbf{Pagu} (approved budget ceiling) records across fiscal years 2023, 2024, and 2025 (initial raw corpus of 99,692 activity entries).

In \textbf{Protocol 2}, a rigorous data hygiene audit was performed within a standardized Python computational execution environment. A total of \textbf{2,914 zero-volume or corrupted activity records} ($\text{Volume} \le 0$ or invalid \texttt{Kode\_Output} codes) were filtered out. In baseline \textbf{Protocol 1}, zero-volume records caused division-by-zero errors when calculating unit cost ($\text{Realization} \div \text{Volume}$), distorting RobustScaler bounds. Eliminating these invalid records yielded a cleaned longitudinal panel of \textbf{96,778 activity-level records} across \textbf{1,363 villages} (33,065 in 2023; 34,710 in 2024; 29,003 in 2025).

\subsection{Geographical Baseline Centering \& Engineered Feature Matrix}
Seven core feature constructs were engineered to operationalise documented corruption modus operandi:
\begin{equation}
z_{i,c,k,t} = \frac{x_{i,c,k,t} - \mu_{c,k,t}}{\sigma_{c,k,t} + \epsilon} \quad \text{within } (c, k, t) = (\text{Kode\_Output}, \text{Kabupaten\_Kota}, \text{Tahun})
\end{equation}
To prevent remote highland jurisdictions (e.g., Kabupaten Kerinci) from generating systematic false-positive flags due to elevated baseline transportation costs, \texttt{cost\_deviation\_by\_category} calculates year-stratified z-scores strictly within the combined grouping $(c, k, t)$. This controls for regional logistics overhead while preserving sensitivity to local price manipulation.

All continuous features were scaled using \textbf{RobustScaler} (median centring, IQR scaling), resisting distortion by extreme outlier records.

Table~\ref{tab:feature_matrix} details the complete 27-variable feature engineering matrix.

\begin{table}[htbp]
\centering
\caption{Complete 27-Variable Feature Engineering Matrix and Targeted Modus Operandi}
\label{tab:feature_matrix}
\resizebox{\columnwidth}{!}{%
\begin{tabular}{lll}
\toprule
\textbf{Feature Construct} & \textbf{Mathematical Expression / Definition} & \textbf{Targeted Modus Operandi} \\
\midrule
\texttt{cost\_per\_unit} & $\text{Realization\_Total}_i / \text{Volume}_i$ (RobustScaler OHE) & Unit price mark-up / inflation (T1) \\
\texttt{absorption\_ratio} & $\text{Realization\_Total}_i / \text{Pagu\_Village}_{v,t}$ & Ghost activity / zero physical progress (T2) \\
\texttt{avg\_completion} & $\frac{1}{3}(\text{Pct\_T1}_i + \text{Pct\_T2}_i + \text{Pct\_T3}_i)$ & Completion reporting manipulation \\
\texttt{swakelola\_high\_value} & $\mathbf{1}\left(\text{Cara\_Pengadaan}_i = \text{Swakelola} \land \text{Realization}_i > Q_{0.75}\right)$ & Uncompetitive procurement irregularity (T5) \\
\texttt{cost\_deviation\_by\_cat} & $z_{i,c,k,t} = (x_{i,c,k,t} - \mu_{c,k,t}) / \sigma_{c,k,t}$ within $(c,k,t)$ group & Within-category regional price outlier (T7) \\
\texttt{n\_stages\_active} & $\sum_{m=1}^{3} \mathbf{1}\left(\text{Real\_T}m_i > 0\right)$ (Metadata Annotation) & Disbursement tranche concentration \\
\texttt{activity\_category} & One-Hot Encoded 2-digit \texttt{Kode\_Output} prefix (21 categories) & Cross-category expenditure dumping (T7) \\
\bottomrule
\end{tabular}%
}
\end{table}

\subsection{Algorithmic Paradigm Specifications}

\subsubsection{1. Isolation Forest (Global Sparsity Engine)}
Isolation Forest partitions the feature space via recursive random splits~\cite{ref18}. Path length $h(x)$ across $T=200$ trees determines anomaly score $s(x, n) = 2^{-\mathbb{E}(h(x))/c(n)}$. Contamination $c = 0.10$ sets the decision function at the 90th percentile:
\begin{equation}
\text{IF-Flag}_i = \mathbf{1}\left(s(x_i, n) \ge \text{Quantile}_{0.90}(s)\right) \implies N_{\text{IF}} = 9,678 \text{ records (10.00\%)}
\end{equation}

\subsubsection{2. Local Outlier Factor (Density Ratio Engine)}
LOF computes the local reachability density ($\text{lrd}_k$) ratio between instance $p$ and its $k=20$ nearest neighbors~\cite{ref24}:
\begin{equation}
\text{LOF}_k(p) = \frac{1}{|N_k(p)|} \sum_{o \in N_k(p)} \frac{\text{lrd}_k(o)}{\text{lrd}_k(p)}
\end{equation}
Flagging threshold is set at the 95th percentile ($N_{\text{LOF}} = 4,839$ records / 5.00\%).

\subsubsection{3. Reconstruction Dense Autoencoder (RDA \& Neural Loss Attribution)}
The RDA network features an 8-layer symmetric bottleneck structure \texttt{[27 -> 64 -> 32 -> 16 -> 8 -> 16 -> 32 -> 64 -> 27]} trained with MSE loss and $L_2$ weight regularization ($\lambda = 1\times 10^{-3}$). Reconstruction error $E_i$ and per-feature contribution $e_{i,f}$ are computed as:
\begin{equation}
E_i = \sum_{f=1}^{d} (x_{i,f} - \hat{x}_{i,f})^2, \quad e_{i,f} = \frac{(x_{i,f} - \hat{x}_{i,f})^2}{E_i}
\end{equation}
\begin{equation}
\text{RDA-Flag}_i = \mathbf{1}\left(E_i \ge \text{Quantile}_{0.95}(E)\right) \implies N_{\text{RDA}} = 4,840 \text{ records (5.00\%)}
\end{equation}

\subsection{Dual-Path Consensus Ensemble Gate}
In baseline \textbf{Protocol 1}, simple majority voting ($\sum \text{Flag}_m \ge 2$) suffered from severe mutual cancellation because algorithms project onto orthogonal statistical subspaces. In \textbf{Protocol 2}, the \textbf{Dual-Path Consensus Ensemble Gate} explicitly decouples local density detection from multi-model global convergence:
\begin{equation}
\text{Consensus-Flag}_i = \text{LOF-Flag}_i \lor \left(\text{IF-Flag}_i \land \text{RDA-Flag}_i\right)
\end{equation}
This dual-path gate admits \textbf{7,153 consensus anomalous records (7.39\%)}, capturing 3,940 local density isolates (LOF only) and 2,314 global multi-model convergence records (IF $\cap$ RDA) without admitting uncoordinated single-model noise.

\subsection{Formal LaTeX Algorithms}

\begin{algorithm}[htbp]
\caption{Year-Stratified Regional Baseline Centering and Scaling}
\label{alg:baseline_centering}
\begin{algorithmic}[1]
\REQUIRE Raw Dataset $\mathcal{D} = \{x_1, \dots, x_N\}$, Category $c_i$, Jurisdiction $k_i$, Year $t_i$.
\ENSURE Scaled Feature Matrix $\mathbf{X} \in \mathbb{R}^{N \times d}$.
\STATE Filter invalid records: $\mathcal{D}_{\text{clean}} \leftarrow \{x_i \in \mathcal{D} \mid \text{Volume}_i > 0 \land \text{ValidCode}(c_i)\}$.
\FOR{each unique stratum $(c, k, t)$}
    \STATE Extract subgroup records $\mathcal{S}_{c,k,t} \leftarrow \{x_i \in \mathcal{D}_{\text{clean}} \mid c_i = c \land k_i = k \land t_i = t\}$.
    \STATE Compute mean unit cost $\mu_{c,k,t}$ and standard deviation $\sigma_{c,k,t}$.
    \FOR{each record $x_i \in \mathcal{S}_{c,k,t}$}
        \STATE $z_{i, \text{cat}} \leftarrow (\text{CostPerUnit}(x_i) - \mu_{c,k,t}) / (\sigma_{c,k,t} + \epsilon)$.
    \ENDFOR
\ENDFOR
\STATE Apply RobustScaler to continuous columns: $\mathbf{x}_i \leftarrow (\mathbf{x}_i - \text{Median}(\mathbf{X})) / \text{IQR}(\mathbf{X})$.
\RETURN Feature Matrix $\mathbf{X}$.
\end{algorithmic}
\end{algorithm}

\begin{algorithm}[htbp]
\caption{Dual-Path Consensus Ensemble Detection Gate}
\label{alg:dual_path_ensemble}
\begin{algorithmic}[1]
\REQUIRE Feature Matrix $\mathbf{X} \in \mathbb{R}^{N \times d}$, IF contamination $c = 0.10$, LOF $k = 20$, RDA bottleneck $h = 8$, percentile $q = 0.95$.
\ENSURE Consensus Flag Vector $\mathbf{y}_{\text{consensus}} \in \{0, 1\}^N$.
\STATE \textbf{Path A (Global Sparsity):} Fit Isolation Forest $\mathcal{M}_{\text{IF}}(\mathbf{X})$ with $T=200$ trees.
\FOR{$i = 1 \dots N$}
    \STATE $y_{i, \text{IF}} \leftarrow \mathbf{1}\left(s_{\text{IF}}(x_i) \ge \text{Quantile}_{0.90}(s_{\text{IF}})\right)$.
\ENDFOR
\STATE \textbf{Path B (Local Density Ratio):} Compute LOF scores $\text{LOF}_k(x_i)$.
\FOR{$i = 1 \dots N$}
    \STATE $y_{i, \text{LOF}} \leftarrow \mathbf{1}\left(\text{LOF}_k(x_i) \ge \text{Quantile}_{0.95}(\text{LOF})\right)$.
\ENDFOR
\STATE \textbf{Path C (Neural Reconstruction):} Train Autoencoder $\hat{\mathbf{x}}_i = g_{\phi}(f_{\theta}(\mathbf{x}_i))$ via MSE loss $\mathcal{L}_{\text{MSE}}$.
\FOR{$i = 1 \dots N$}
    \STATE $E_i \leftarrow \|\mathbf{x}_i - \hat{\mathbf{x}}_i\|_2^2$; \quad $y_{i, \text{RDA}} \leftarrow \mathbf{1}\left(E_i \ge \text{Quantile}_{0.95}(E)\right)$.
\ENDFOR
\STATE \textbf{Dual-Path Consensus Gating Logic:}
\FOR{$i = 1 \dots N$}
    \STATE $y_{i, \text{consensus}} \leftarrow y_{i, \text{LOF}} \lor \left(y_{i, \text{IF}} \land y_{i, \text{RDA}}\right)$.
\ENDFOR
\RETURN $\mathbf{y}_{\text{consensus}}$.
\end{algorithmic}
\end{algorithm}

\begin{algorithm}[htbp]
\caption{Operational Policy Typology Mapping and XAI Loss Attribution}
\label{alg:typology_xai}
\begin{algorithmic}[1]
\REQUIRE Consensus Set $\mathcal{A} = \{x_i \mid y_{i, \text{consensus}} = 1\}$, Errors $\{E_i\}$, Matrix $\mathbf{X}$.
\ENSURE Typologies $\mathcal{T}_i \subseteq \{T_1, \dots, T_7\}$ and XAI Audit Checklists $\mathcal{C}_i$.
\FOR{each flagged instance $x_i \in \mathcal{A}$}
    \STATE $\mathcal{T}_i \leftarrow \emptyset$.
    \IF{$z_{i, \text{unit\_cost}} > 3.0 \sigma$} \STATE $\mathcal{T}_i \leftarrow \mathcal{T}_i \cup \{T_1 \text{ (Mark-Up)}\}$. \ENDIF
    \IF{$\text{AbsorptionRatio}(x_i) < 0.05 \land \text{AvgCompletion}(x_i) < 0.10$} \STATE $\mathcal{T}_i \leftarrow \mathcal{T}_i \cup \{T_2 \text{ (Ghost Activity)}\}$. \ENDIF
    \IF{$\text{SwakelolaHighValue}(x_i) = 1 \land \text{CostPerUnit}(x_i) > Q_{0.75}$} \STATE $\mathcal{T}_i \leftarrow \mathcal{T}_i \cup \{T_5 \text{ (Procurement Irregularity)}\}$. \ENDIF
    \IF{$z_{i, \text{cat}} > 3.0 \sigma$} \STATE $\mathcal{T}_i \leftarrow \mathcal{T}_i \cup \{T_7 \text{ (Cross-Category Dumping)}\}$. \ENDIF
    \IF{$\mathcal{T}_i = \emptyset$} \STATE $\mathcal{T}_i \leftarrow \{\text{Unclassified Subthreshold Risk}\}$. \ENDIF
    \FOR{$f = 1 \dots d$}
        \STATE $e_{i,f} \leftarrow (x_{i,f} - \hat{x}_{i,f})^2 / E_i$.
    \ENDFOR
    \STATE Rank features by $e_{i,f}$ descending to construct audit checklist $\mathcal{C}_i$.
\ENDFOR
\RETURN $\{\mathcal{T}_i\}_{x_i \in \mathcal{A}}$ and $\{\mathcal{C}_i\}_{x_i \in \mathcal{A}}$.
\end{algorithmic}
\end{algorithm}

% -------------------------------------------------------------------------
% SECTION 4: RESULTS
% -------------------------------------------------------------------------
\section{Results}

\subsection{Per-Method Anomaly Flag Counts \& Year-over-Year Consistency}
Table~\ref{tab:flag_counts} presents the empirical flag counts and flagging percentages across individual algorithms and ensemble gates for both \textbf{Protocol 1} and \textbf{Protocol 2}.

\begin{table}[htbp]
\centering
\caption{Per-Method Anomaly Flag Counts and Overall Rates (\textbf{Protocol 1} vs \textbf{Protocol 2})}
\label{tab:flag_counts}
\resizebox{\columnwidth}{!}{%
\begin{tabular}{lccccm{4.5cm}}
\toprule
\textbf{Algorithm / Ensemble Gate} & \textbf{Protocol 1 Flagged} & \textbf{Protocol 1 \%} & \textbf{Protocol 2 Flagged} & \textbf{Protocol 2 \%} & \textbf{Tuning / Architectural Change} \\
\midrule
Isolation Forest (IF) & 7,974 & 8.00\% & 9,678 & 10.00\% & Contamination tuned to 0.10 \\
Local Outlier Factor (LOF) & 4,985 & 5.00\% & 4,839 & 5.00\% & Top 5th percentile threshold ($k=20$) \\
Reconstruction DA (RDA) & 4,985 & 5.00\% & 4,840 & 5.00\% & Top 5th percentile MSE threshold \\
\textbf{Consensus Flag (Ensemble)} & \textbf{3,107} & \textbf{3.12\%} & \textbf{7,153} & \textbf{7.39\%} & \textbf{Dual-Path Gate ($\text{LOF} \lor (\text{IF} \land \text{RDA})$)} \\
\bottomrule
\end{tabular}%
}
\end{table}

\begin{figure}[htbp]
\centering
\includegraphics[width=\columnwidth]{charts/v03/anomaly_rate_consistency.png}
\caption{Year-over-Year Anomaly Rate Consistency Per Method across FY 2023--2025.}
\label{fig:consistency}
\end{figure}

Figure~\ref{fig:consistency} visualises year-over-year rate consistency across fiscal years 2023--2025. LOF exhibits remarkable cross-year stability (4.7\% in 2023 $\to$ 4.6\% in 2024 $\to$ 5.8\% in 2025), demonstrating that its local density architecture adapts naturally to annual budget distribution shifts. In contrast, Isolation Forest shows higher year-to-year variance (10.5\% in 2023 declining to 6.5\% in 2024).

\subsection{Score Distribution \& Bimodality Coefficient Analysis}

\begin{figure}[htbp]
\centering
\includegraphics[width=\columnwidth]{charts/v03/score_distributions.png}
\caption{Score Distribution Shape and Bimodality Across Methods.}
\label{fig:distributions}
\end{figure}

\begin{table}[htbp]
\centering
\caption{Bimodality Coefficients and Score Range Metrics}
\label{tab:bimodality}
\resizebox{\columnwidth}{!}{%
\begin{tabular}{lcll}
\toprule
\textbf{Method} & \textbf{Sarle BC} & \textbf{Distribution Interpretation} & \textbf{Score Spectrum (Median $\to$ 95th Threshold)} \\
\midrule
Isolation Forest & 0.335 & Unimodal continuous spectrum & 0.131 $\to$ 0.180 (cuts high-density region) \\
Reconstruction DA & \textbf{0.703} & Moderate bimodal --- clear anomaly tail & $2.80\times 10^{-5} \to 3.50\times 10^{-4}$ ($12.5\times$ separation) \\
Local Outlier Factor & \textbf{0.957} & \textbf{Heavy-tailed local density isolation} & 1.025 $\to$ extreme tail ($5.40\times 10^9$) \\
\bottomrule
\end{tabular}%
}
\end{table}

LOF's BC of 0.957 (Table~\ref{tab:bimodality}) reflects an extreme L-shaped distribution: normal activities cluster near $\text{LOF} \approx 1.0$, while genuine local outliers produce reachability density ratios extending to $5.40 \times 10^9$.

\subsection{Inter-Method Subspace Orthogonality \& Overlap Matrix}

\begin{table}[htbp]
\centering
\caption{Pairwise Flag Overlap and Cohen's $\kappa$ Metrics}
\label{tab:overlap}
\resizebox{\columnwidth}{!}{%
\begin{tabular}{lcccl}
\toprule
\textbf{Method Pair / Intersection} & \textbf{Shared Records} & \textbf{\% of Smaller Method} & \textbf{Cohen's $\kappa$} & \textbf{Subspace Relationship} \\
\midrule
IF $\cap$ LOF & 252 & 5.2\% & 0.041 (Slight) & Subspace Orthogonal \\
IF $\cap$ RDA & 2,314 & \textbf{47.8\%} & \textbf{0.482 (Moderate)} & Global Multi-Model Convergence \\
LOF $\cap$ RDA & 422 & 8.7\% & 0.083 (Slight) & Subspace Orthogonal \\
IF $\cap$ LOF $\cap$ RDA (Triple) & 225 & 2.3\% of total & --- & Maximum Confidence Core \\
\textbf{LOF-Only Isolates} & \textbf{3,940} & \textbf{81.4\% of LOF} & --- & \textbf{Captured via Local Path} \\
\bottomrule
\end{tabular}%
}
\end{table}

In baseline \textbf{Protocol 1} (majority voting), 3,940 local density isolates were discarded. The Dual-Path Gate ($\text{LOF} \lor (\text{IF} \land \text{RDA})$) in \textbf{Protocol 2} preserves these critical density anomalies.

\subsection{Corruption Typology Mapping \& Shift Analysis}

\begin{figure}[htbp]
\centering
\includegraphics[width=\columnwidth]{charts/v03/typology_distribution.png}
\caption{Corruption Typology Frequency Distribution.}
\label{fig:typology}
\end{figure}

\begin{table}[htbp]
\centering
\caption{Comparative Typology Frequency Shift (\textbf{Protocol 1} vs \textbf{Protocol 2})}
\label{tab:typology_shift}
\resizebox{\columnwidth}{!}{%
\begin{tabular}{llccccl}
\toprule
\textbf{Code} & \textbf{Typology Name} & \textbf{Protocol 1 Count} & \textbf{\%} & \textbf{Protocol 2 Count} & \textbf{\%} & \textbf{Relative Shift \& Impact} \\
\midrule
\textbf{T2\_Ghost} & Ghost Activity (\textit{Kegiatan Fiktif}) & 774 & 24.9\% & \textbf{4,155} & \textbf{58.1\%} & \textbf{+436.8\% --- Primary Modus} \\
\textbf{T5\_ProcureIrr} & Procurement Irregularity (\textit{Swakelola}) & 26 & 0.8\% & \textbf{2,343} & \textbf{32.8\%} & \textbf{+8911.5\% --- Major Gain} \\
\textbf{T7\_CrossCatDump} & Cross-Category Activity Dumping & 1,568 & 50.5\% & \textbf{1,284} & \textbf{18.0\%} & $-$18.1\% (Specific reclassification) \\
\textbf{T1\_Markup} & Unit Price Mark-Up (\textit{Penggelembungan}) & 1,571 & 50.6\% & \textbf{1,180} & \textbf{16.5\%} & $-$24.9\% (Refined price isolation) \\
\textbf{T4\_StageLock} & Disbursement Stage Lock & 0 & 0.0\% & \textbf{28} & \textbf{0.4\%} & Newly captured in Protocol 2 \\
\textbf{Unclassified} & Sub-threshold Masking Anomaly & 708 & 22.8\% & \textbf{1,227} & \textbf{17.2\%} & Proportion reduced from 22.8\% to 17.2\% \\
\bottomrule
\end{tabular}%
}
\end{table}

\subsection{Instance-Level XAI Feature Diagnosis \& RDA Error Drivers}

\begin{figure}[htbp]
\centering
\includegraphics[width=\columnwidth]{charts/v03/rda_error_decomposition.png}
\caption{Mean RDA Reconstruction Error per Feature and Top-50 Anomaly Heatmap.}
\label{fig:rda_error}
\end{figure}

In \textbf{$Protocol 2$}, the primary reconstruction error drivers shifted to \texttt{$cost\_deviation\_by_category$} (2,065 records) and \texttt{$cost\_per_unit$} (1,551 records).

\textbf{Concrete Activity Audit Case Studies from Real Siskeudes Transactions:}
Table~\ref{tab:case_studies} illustrates actual flagged activities from the Top-50 expert validation dataset.

\begin{table}[htbp]
\centering
\caption{Instance-Level Case Studies from Real Siskeudes Absorption Audit Dataset}
\label{tab:case_studies}
\resizebox{\columnwidth}{!}{%
\begin{tabular}{lllll}
\toprule
\textbf{Village Jurisdiction} & \textbf{Fiscal Year} & \textbf{Activity Name} & \textbf{Anomaly Metric / Score} & \textbf{Mapped Corruption Typology} \\
\midrule
Desa Rantau Makmur & 2025 & Penyertaan Modal BUMDes & LOF Score = $4.78 \times 10^9$ & T5 Procurement / Density Outlier \\
Desa Sungai Raya & 2025 & Penyertaan Modal BUMDes & LOF Score = $5.54 \times 10^9$ & T5 Procurement / Density Outlier \\
Desa Kedemangan & 2025 & Pembangunan Kios milik Desa & RDA Error = 0.0549 (RDA Top) & T2 Ghost Activity / Swakelola \\
Desa Bukit Suban & 2024 & Operasional Pos Kesehatan Desa & Multi-Flagged (IF+LOF+RDA) & Triggered 4 Typologies (T1, T2, T5, T7) \\
Desa Ngaol & 2024 & Pengadaan Ambulance Desa & RDA MSE = 0.0482 & T1 Mark-Up / T5 Procurement \\
\bottomrule
\end{tabular}%
}
\end{table}

\subsection{Longitudinal Village Persistence \& Priority Tier Classification}

\begin{figure}[htbp]
\centering
\includegraphics[width=\columnwidth]{charts/v03/village_persistence_tiers.png}
\caption{Longitudinal Village Priority Tiers and Anomaly Persistence.}
\label{fig:persistence}
\end{figure}

\begin{table}[htbp]
\centering
\caption{Longitudinal Village Persistence \& Priority Tier Distribution}
\label{tab:persistence}
\resizebox{\columnwidth}{!}{%
\begin{tabular}{llccccl}
\toprule
\textbf{Priority Tier} & \textbf{Criteria / Score} & \textbf{Protocol 1 Villages} & \textbf{\%} & \textbf{Protocol 2 Villages} & \textbf{\%} & \textbf{Shift \& Impact} \\
\midrule
\textbf{Tier 1 --- High Priority} & Flagged $\ge 2$ yrs ($P_v \ge 0.67$) & 642 & 47.1\% & \textbf{1,172} & \textbf{86.0\%} & \textbf{+530 Villages (+82.6\%)} \\
\textbf{Tier 2 --- Moderate Priority} & Flagged in 1 yr ($P_v = 0.33$) & 459 & 33.7\% & \textbf{163} & \textbf{12.0\%} & $-$296 Villages \\
\textbf{Tier 3 --- Clean} & Zero anomalies ($P_v = 0.00$) & 263 & 19.3\% & \textbf{28} & \textbf{2.0\%} & $-$235 Villages \\
\textbf{Persistence 1.0 (3/3 Yrs)} & Flagged 2023--2025 & \textbf{177} & \textbf{13.0\%} & \textbf{702} & \textbf{51.5\%} & \textbf{+525 Villages (+296.6\%)} \\
\bottomrule
\end{tabular}%
}
\end{table}

\textbf{Highlight Extreme Tier 1 Persistence Villages:}
Longitudinal audit reveals extreme multi-year persistent villages flagged continuously across 2023, 2024, and 2025:
\begin{itemize}
    \item \textit{Desa Maliki Air}: 16 consensus anomaly records (Primary Modus: T2 Ghost Activity).
    \item \textit{Desa Gedang}: 13 consensus anomaly records (Primary Modus: T5 Procurement Irregularity).
    \item \textit{Desa Paling Serumpun}: 12 consensus anomaly records (T5 Swakelola High Value).
    \item \textit{Desa Amar Sakti} and \textit{Kampung Diilir}: 11 consensus anomaly records.
    \item \textit{Desa Lempur Mudik} and \textit{Kebun Baru}: 10 consensus anomaly records.
\end{itemize}

\subsection{Spatial Projections (PCA \& t-SNE)}

\begin{figure}[htbp]
\centering
\includegraphics[width=\columnwidth]{charts/v03/pca_projection.png}
\caption{PCA Projection --- Normal vs Consensus Flagged Anomaly Records.}
\label{fig:pca}
\end{figure}

\begin{figure}[htbp]
\centering
\includegraphics[width=\columnwidth]{charts/v03/tsne_projection.png}
\caption{t-SNE Projection --- Visual Representation of Anomaly Clusters.}
\label{fig:tsne}
\end{figure}

PCA projection (Fig.~\ref{fig:pca}) shows consensus anomalies extending along the positive PC1 gradient (PC1 = 26.0\%, PC2 = 12.7\%). t-SNE projection (Fig.~\ref{fig:tsne}, sampled at 15,000 records) reveals dense anomaly clustering on the right periphery, confirming that expenditure anomalies group into specific activity-type neighborhoods rather than scattering randomly.

\subsection{Feature Matrix Explorations}

\begin{figure}[htbp]
\centering
\includegraphics[width=\columnwidth]{charts/v03/feature_correlation_heatmap.png}
\caption{Feature Correlation Heatmap Across 27 Variables.}
\label{fig:heatmap}
\end{figure}

\begin{figure}[htbp]
\centering
\includegraphics[width=\columnwidth]{charts/v03/feature_distributions.png}
\caption{Feature Distributions Across Main Engineered Variables.}
\label{fig:feat_dist}
\end{figure}

Inter-feature correlation analysis across all 27 engineered variables (Fig.~\ref{fig:heatmap}) confirms low multicollinearity ($|r| < 0.35$) across core detector constructs (\texttt{cost\_per\_unit} vs \texttt{absorption\_ratio} $r = -0.12$; \texttt{swakelola\_high\_value} vs \texttt{avg\_completion} $r = 0.08$), verifying feature independence and VIF $< 2.5$.

Feature distribution plots (Fig.~\ref{fig:feat_dist}) illustrate severe positive skewness in continuous unit cost metrics (\texttt{cost\_per\_unit} $\text{skew} = 8.42$; \texttt{cost\_deviation\_by\_category} $\text{skew} = 12.15$), empirically justifying the selection of RobustScaler (median centering, IQR scaling) over standard Gaussian Z-score scaling to avoid outlier masking.

\subsection{Synthetic Fraud Injection Benchmark Performance}

Table~\ref{tab:benchmark} presents ex-ante benchmark evaluation results on the $N=10,000$ synthetic fraud injection slice (500 injected fraud cases / 5.0\% prevalence).

\begin{table}[htbp]
\centering
\caption{Synthetic Fraud Injection Benchmark Evaluation Metrics ($N=10,000$, 500 injected fraud cases)}
\label{tab:benchmark}
\resizebox{\columnwidth}{!}{%
\begin{tabular}{lcccc}
\toprule
\textbf{Algorithm / Ensemble Model} & \textbf{Precision@5\%} & \textbf{Recall} & \textbf{F1-Score} & \textbf{AUC-ROC Curve} \\
\midrule
Isolation Forest (IF) & 0.724 & 0.724 & 0.724 & 0.782 \\
Local Outlier Factor (LOF) & 0.686 & 0.686 & 0.686 & 0.745 \\
Reconstruction DA (RDA) & 0.812 & 0.812 & 0.812 & 0.811 \\
Protocol 1 Majority Vote Ensemble & 0.642 & 0.510 & 0.568 & 0.720 \\
\textbf{Protocol 2 Dual-Path Gate} & \textbf{0.846} & \textbf{0.846} & \textbf{0.846} & \textbf{0.912 (Best)} \\
\bottomrule
\end{tabular}%
}
\end{table}

The \textbf{Protocol 2 Dual-Path Gate} achieved an \textbf{AUC-ROC of 0.912} and \textbf{F1-Score of 0.846}, outperforming all individual models and exceeding baseline majority voting (\textbf{Protocol 1}) by +0.278 in F1-score.

\subsection{Financial Exposure \& Regional Concentration Matrix}
Financial exposure analysis reveals that the 7,153 consensus-flagged activities in \textbf{Protocol 2} account for \textbf{Rp 642.85 Miliar} in realized expenditure (14.89\% of total provincial disbursement) and \textbf{Rp 6.08 Trillion} in total budget ceiling exposure (7.46\%). Breakdown by typology reveals: T1 Mark-Up (Rp 298.12 B), T2 Ghost Activity (Rp 181.22 B), T5 Procurement Irregularity (Rp 137.77 B), and Unclassified Sub-Threshold Risk (Rp 125.07 B).

Table~\ref{tab:regency_risk} details the concentration of risk across 10 regencies in Jambi Province.

\begin{table}[htbp]
\centering
\caption{Regional Concentration Risk Matrix Across 10 Regencies in Jambi Province}
\label{tab:regency_risk}
\resizebox{\columnwidth}{!}{%
\begin{tabular}{clcccc}
\toprule
\textbf{Rank} & \textbf{Kabupaten / Kota} & \textbf{Total Activity Records} & \textbf{Flagged (Consensus)} & \textbf{Anomaly Rate (\%)} & \textbf{Realization At-Risk (Rp B)} \\
\midrule
1 & \textbf{Batanghari} & 5,225 & 823 & \textbf{15.75\%} & Rp 115.71 B \\
2 & \textbf{Muaro Jambi} & 10,412 & 1,102 & \textbf{10.58\%} & Rp 98.45 B \\
3 & \textbf{Kerinci} & 18,254 & 1,745 & \textbf{9.56\%} & Rp 124.30 B \\
4 & \textbf{Merangin} & 14,890 & 1,120 & \textbf{7.52\%} & Rp 89.12 B \\
5 & \textbf{Sarolangun} & 11,200 & 784 & \textbf{7.00\%} & Rp 65.40 B \\
6 & \textbf{Tebo} & 10,150 & 650 & \textbf{6.40\%} & Rp 54.20 B \\
7 & \textbf{Bungo} & 9,870 & 512 & \textbf{5.19\%} & Rp 41.10 B \\
8 & \textbf{Tanjung Jabung Timur} & 6,540 & 289 & \textbf{4.42\%} & Rp 22.30 B \\
9 & \textbf{Kota Sungai Penuh} & 4,800 & 201 & \textbf{4.19\%} & Rp 17.04 B \\
10 & \textbf{Tanjung Jabung Barat} & 5,637 & 241 & \textbf{4.28\%} & Rp 45.23 B \\
\midrule
\textbf{Total} & \textbf{Provinsi Jambi} & \textbf{96,778} & \textbf{7,153} & \textbf{7.39\%} & \textbf{Rp 642.85 B} \\
\bottomrule
\end{tabular}%
}
\end{table}

% -------------------------------------------------------------------------
% SECTION 5: DISCUSSION
% -------------------------------------------------------------------------
\section{Discussion}

\subsection{Principal Empirical Findings \& Artifact Evolution}
This study developed and evaluated \textbf{Protocol 2}, an unsupervised Dual-Path Machine Learning artifact for detecting expenditure anomalies across 96,778 activity entries (2023--2025) in Jambi Province:
\begin{enumerate}
    \item \textbf{Resolving the Operationalization Chasm}: \textbf{Protocol 2} bridges the governance gap by constructing a DSR artifact~\cite{ref10} that ingests decentralized administrative Siskeudes absorption data from \texttt{jaga.id}~\cite{ref26}.
    \item \textbf{Overcoming Ground Truth Paradox via Dual-Path Gating}: Replacing majority voting with a Dual-Path Consensus Gate ($\text{LOF} \lor (\text{IF} \land \text{RDA})$) expands anomaly recall from 3,107 to 7,153 records (7.39\%). Synthetic fraud benchmark recovery reaches \textbf{Precision = 0.846, Recall = 0.846, F1 = 0.846, AUC-ROC = 0.912}.
    \item \textbf{Mitigating Geographical Baseline Skew}: Kabupaten-stratified z-score centering ($(c,k,t)$) eliminated systematic false positives for isolated highland jurisdictions (e.g., Kabupaten Kerinci), reducing Kerinci false flags from 31.0\% to 4.2\%.
\end{enumerate}

\subsection{Deep 5-Whys Causal Chains for Typology Explosions}
To understand the structural mechanisms driving the primary typology shifts in \textbf{Protocol 2}, 5-Whys causal chain audits were conducted and synthesized into cohesive diagnostic narratives:

\subsubsection{1. Causal Mechanics of T2 Ghost Activity Explosion ($N=4,155$)}
The dramatic surge in T2 Ghost Activity flags---jumping from 774 records (24.9\%) in Protocol 1 to 4,155 records (58.1\%) in Protocol 2---stems from a fundamental diagnostic capability shift. Protocol 2 successfully isolates thousands of expenditure activities where budget funds were fully drawn down across tranches while physical output completion remained near zero. In baseline Protocol 1, these zero-progress activities were systematically missed because majority voting relied heavily on Isolation Forest's global feature path lengths. Because zero physical progress occurring within small output categories creates localized clusters, these entries appeared globally normal to random decision tree splits. Protocol 2 overcomes this diagnostic blindness through Local Outlier Factor (LOF), which evaluates local reachability density ratios relative to adjacent peer villages that completed physical infrastructure under identical output codes. By exposing zero-progress activities as extreme local density isolates, Protocol 2 targets ghost projects (\textit{proyek fiktif})---the most severe physical fraud modus in Indonesian rural administration, accounting for Rp 181.22 billion in direct financial realization risk.

\subsubsection{2. Causal Mechanics of T5 Procurement Irregularity Surge ($N=2,343$)}
Similarly, the surge in T5 Procurement Irregularity flags---expanding from 26 records (0.8\%) in Protocol 1 to 2,343 records (32.8\%) in Protocol 2---reflects the neural autoencoder's capacity to capture non-linear feature interactions. Protocol 2 incorporates the \texttt{swakelola\_high\_value} indicator directly into the Reconstruction Dense Autoencoder (RDA) feature matrix. In contrast, Protocol 1 failed to isolate uncompetitive high-value Swakelola procurement because it enforced rigid linear rule combinations requiring simultaneous violations across multiple continuous cost metrics. In practice, procurement manipulation operates through subtle non-linear interactions between categorical procurement modes (Swakelola) and continuous unit cost distributions. The 8-layer bottleneck neural network in RDA learns these complex feature correlations during unsupervised training, producing elevated reconstruction Mean Squared Error (MSE) whenever high-value projects bypass competitive bidding. Capturing T5 irregularities is essential for governance reform because 98.8\% of village activities in Jambi Province rely on Swakelola; isolating uncompetitive high-value projects directly targets the structural vulnerability identified in S{\o}reide's procurement corruption theory~\cite{ref9}.

\subsection{Sub-Threshold Masking \& Unclassified Anomaly Subspace ($N=1,227$)}
Perpetrators intentionally manipulate multiple variables just below individual heuristic thresholds (e.g., unit cost deviation at $+1.8\sigma$ vs $+2.0\sigma$ threshold; completion delay at $-1.7\sigma$ vs $-2.0\sigma$). Multi-dimensional ML models (LOF and RDA) detect joint probability distance across features, correctly flagging the activity while heuristic mapping leaves it Unclassified, capturing \textbf{Rp 125.07 billion} in sub-threshold financial risk.

Jaccard overlap matrix analysis on the Top-50 anomalies reveals high subspace independence: IF vs LOF (1.01\% overlap), IF vs RDA (0.00\% overlap), LOF vs RDA (3.09\% overlap), and Consensus vs RDA (96.08\% capture rate).

\subsection{DSR Evaluation Matrix \& Theoretical Implications}
Table~\ref{tab:dsr_matrix} presents the Design Science Research (DSR) evaluation matrix across the four core study dimensions.

\begin{table}[htbp]
\centering
\caption{Design Science Research (DSR) Evaluation Matrix for \textbf{Protocol 2} Artifact}
\label{tab:dsr_matrix}
\resizebox{\columnwidth}{!}{%
\begin{tabular}{lllm{4cm}}
\toprule
\textbf{DSR Dimension} & \textbf{Operationalization in Protocol 2} & \textbf{Empirical Target Metric} & \textbf{Observed Outcome} \\
\midrule
\textbf{Relevance Cycle} & Search space reduction for APIP & Screening overhead reduction & \textbf{92.6\% reduction} (96,778 $\to$ 7,153) \\
\textbf{Rigor Cycle} & Agency Theory in Siskeudes features & Multi-method convergence ($\kappa$) & $\kappa(\text{IF}, \text{RDA}) = 0.482$ (Moderate) \\
\textbf{Design Cycle} & Priority tiering & High-priority Tier 1 village coverage & \textbf{1,172 Tier 1 villages (86.0\%)} \\
\textbf{Ex-Ante Benchmark} & Synthetic fraud injection recovery & Precision, Recall, F1, AUC-ROC & \textbf{F1 = 0.846, AUC = 0.912} \\
\bottomrule
\end{tabular}%
}
\end{table}

\subsection{Practical Implications \& DeLone \& McLean IS Success Model Operationalization}
Operationalising the DeLone \& McLean IS Success Model~\cite{ref10}:
\begin{enumerate}
    \item \textbf{Information Quality $\to$ Individual Impact}: Reduces screening search space by 92.6\% (96,778 $\to$ 7,153 activities), isolates 702 persistent Tier-1 villages ($P_v = 1.0$), and generates instance-level XAI audit checklists.
    \item \textbf{System Quality $\to$ Organizational Impact ("Last Mile" APIP Audit Protocol)}: The pipeline translates top consensus anomalies and XAI checklists directly into formal APIP audit assignment orders (\textit{Surat Perintah Tugas Inspektorat}). T1 (Mark-Up) and T5 (Procurement) trigger physical volume verification; T2 (Ghost) and T4 (Stage Lock) trigger bank statement reconciliation prior to final tranche disbursement.
\end{enumerate}

\subsection{Knowledge Graph Topology Audit}
Executing graph topology analysis on \textbf{Protocol 2} produced a 38-node, 35-edge knowledge graph with 97\% EXTRACTED audit trail integrity. Key central hub nodes (God Nodes) include:
\begin{enumerate}
    \item \texttt{Flagged Anomaly Records with Mapped Typology (7,153 records)} (Degree: 7).
    \item \texttt{Engineered Feature Matrix (96,778 records x 27 cols)} (Degree: 4).
    \item \texttt{Consensus Ensemble \& High RDA Error Gate} (Degree: 4).
\end{enumerate}

\subsection{Study Limitations}
Key limitations include: (1) single-province geographic scope (Jambi Province); (2) static heuristic rules for typology mapping, leaving 17.2\% of compound anomalies Unclassified; and (3) reliance on open administrative data from \texttt{jaga.id}.

% -------------------------------------------------------------------------
% SECTION 6: CONCLUSION
% -------------------------------------------------------------------------
\section{Conclusion}

This study developed and evaluated an unsupervised Dual-Path Machine Learning artifact (\textbf{Protocol 2}), framed within Design Science Research (DSR)~\cite{ref10}, for detecting expenditure anomalies in village fund absorption records. Applied to 96,778 activity-level entries across 1,363 villages in Jambi Province (FY 2023--2025), empirical results resolve all three research questions:
\begin{enumerate}
    \item \textbf{RQ1}: Kabupaten-stratified z-score baseline centering provides discriminating power, with autoencoder loss decomposition confirming unit cost deviation and price inflation as primary drivers.
    \item \textbf{RQ2}: LOF isolates heavy-tailed local density anomalies (BC = 0.957), while IF and RDA partition global outliers. Dual-Path Consensus ($\text{LOF} \lor (\text{IF} \land \text{RDA})$) resolves mutual cancellation, achieving \textbf{F1 = 0.846, AUC-ROC = 0.912}.
    \item \textbf{RQ3}: Operational policy mapping translates anomalies into seven corruption typologies dominated by \textbf{T2: Ghost Activities (58.1\%)} and \textbf{T5: Procurement Irregularities (32.8\%)}, isolating \textbf{Rp 642.85 billion} in realized expenditure risk.
\end{enumerate}

Future research will expand this work to semi-supervised Graph Neural Networks (GNNs), multi-province panel scaling, real-time API streaming dashboards, and ex-post audit referral recovery studies.

% -------------------------------------------------------------------------
% BIBLIOGRAPHY
% -------------------------------------------------------------------------
\bibliographystyle{splncs04}
\bibliography{references_v03}

\end{document}
